% !TEX root = ../../main.tex
\subsection{数据库设计}

FoodShield 使用 SQLite 作为数据存储后端，数据库共设计四张核心表：\texttt{users}（用户表）、\texttt{orders}（订单表）、\texttt{messages}（消息表）和 \texttt{audit\_logs}（审计日志表）。各表字段的设计均遵循"最小必要原则"，避免将超出业务所需的敏感信息明文保存在数据库中。

\subsubsection*{用户表（users）}

用户表存储用户的注册信息与匿名身份标识，字段设计如表~\ref{tab:db_users} 所示。

\begin{table}[H]
\centering
\caption{users 表字段说明}
\label{tab:db_users}
\renewcommand{\arraystretch}{1.5}
\setlength{\tabcolsep}{6pt}
\zihao{-4}
\begin{tabular}{L{3.0cm} L{2.5cm} L{3.5cm} L{4.5cm}}
\toprule
\textbf{字段名} & \textbf{类型} & \textbf{说明} & \textbf{安全设计意义} \\
\midrule
id              & INTEGER       & 自增主键               & 内部索引，不对外暴露 \\
\texttt{username}        & TEXT UNIQUE   & 用户名                 & 用于登录认证 \\
\texttt{pid}             & TEXT UNIQUE   & 匿名身份标识（PID）    & 隐藏真实 id，用于通信和溯源 \\
\texttt{pid\_r}          & TEXT          & PID 生成随机数$r$     & 仅服务端保存，用于条件溯源时重算 PID \\
\texttt{password\_hash}  & TEXT          & 口令哈希值             & SM3(salt$\|$password)，不存明文口令 \\
\texttt{salt}            & TEXT          & 随机盐值               & 防彩虹表攻击，每用户独立随机 \\
\texttt{created\_at}     & TIMESTAMP     & 注册时间               & 用于审计和频率限制参考 \\
\bottomrule
\end{tabular}
\end{table}

\subsubsection*{订单表（orders）}

订单表存储订单生命周期数据，字段设计如表~\ref{tab:db_orders} 所示。

\begin{table}[H]
\centering
\caption{orders 表字段说明}
\label{tab:db_orders}
\renewcommand{\arraystretch}{1.5}
\setlength{\tabcolsep}{6pt}
\zihao{-4}
\begin{tabular}{L{3.2cm} L{2.5cm} L{3.3cm} L{4.5cm}}
\toprule
\textbf{字段名} & \textbf{类型} & \textbf{说明} & \textbf{安全设计意义} \\
\midrule
id              & INTEGER       & 自增主键               & 内部索引 \\
\texttt{order\_id}       & TEXT UNIQUE   & UUID 订单号            & 随机生成，防顺序枚举 \\
\texttt{user\_id}        & INTEGER       & 关联用户内部 id        & 外键引用 users.id，不暴露 PID \\
\texttt{token}           & TEXT          & 订单 Token             & HMAC-SM3 认证凭证，入房校验用 \\
\texttt{token\_timestamp}& TEXT          & Token 生成时间戳       & 用于 Token 有效期校验 \\
\texttt{status}          & TEXT          & 订单状态               & 枚举值：created/taken/delivering/completed/cancelled \\
\texttt{note}            & TEXT          & 用户备注/标签          & 对骑手端和管理员端隐藏，仅用户自查 \\
\texttt{created\_at}     & TIMESTAMP     & 订单创建时间           & 用于历史查询与审计 \\
\bottomrule
\end{tabular}
\end{table}

\subsubsection*{消息表（messages）}

消息表存储通信记录，字段设计如表~\ref{tab:db_messages} 所示。

\begin{table}[H]
\centering
\caption{messages 表字段说明}
\label{tab:db_messages}
\renewcommand{\arraystretch}{1.5}
\setlength{\tabcolsep}{6pt}
\zihao{-4}
\begin{tabular}{L{3.2cm} L{2.5cm} L{3.3cm} L{4.5cm}}
\toprule
\textbf{字段名} & \textbf{类型} & \textbf{说明} & \textbf{安全设计意义} \\
\midrule
id              & INTEGER       & 自增主键               & 内部索引，也用于消息顺序 \\
\texttt{msg\_id}         & TEXT UNIQUE   & UUID 消息号            & 全局唯一，防重复提交 \\
\texttt{order\_id}       & TEXT          & 关联订单号             & 外键引用 orders.order\_id \\
\texttt{sender\_pid}     & TEXT          & 发送方 PID             & 不暴露真实用户名 \\
\texttt{role}            & TEXT          & 发送方角色             & 枚举：user/rider/admin/system \\
\texttt{content}         & TEXT          & 消息密文               & SM4-CBC 加密存储，不存明文 \\
\texttt{message\_hash}   & TEXT          & 消息摘要               & SM3 对明文字段联合计算，用于 Merkle Tree 叶子节点 \\
\texttt{timestamp}       & TEXT          & 消息时间戳             & ISO 8601 格式，参与哈希计算 \\
\texttt{created\_at}     & TIMESTAMP     & 写入时间               & 与 timestamp 一致性校验参考 \\
\bottomrule
\end{tabular}
\end{table}

\subsubsection*{审计日志表（audit\_logs）}

审计日志表存储 Merkle Root 快照和管理员操作记录，字段设计如表~\ref{tab:db_audit} 所示。

\begin{table}[H]
\centering
\caption{audit\_logs 表字段说明}
\label{tab:db_audit}
\renewcommand{\arraystretch}{1.5}
\setlength{\tabcolsep}{6pt}
\zihao{-4}
\begin{tabular}{L{3.0cm} L{2.5cm} L{3.5cm} L{4.5cm}}
\toprule
\textbf{字段名} & \textbf{类型} & \textbf{说明} & \textbf{安全设计意义} \\
\midrule
id              & INTEGER       & 自增主键               & 内部索引 \\
\texttt{order\_id}       & TEXT          & 关联订单号             & 可为 NULL（登录等全局操作） \\
\texttt{action}          & TEXT          & 操作类型               & 枚举：ROOT\_UPDATED / VERIFY\_OK / FAIL 等 \\
\texttt{detail}          & TEXT          & 操作详情（JSON）       & 记录消息数、不匹配哈希列表、管理员用户名等 \\
\texttt{merkle\_root}    & TEXT          & 当次 Merkle Root       & 快照值，供后续比对使用 \\
\texttt{created\_at}     & TIMESTAMP     & 操作时间               & 审计时间线依据 \\
\bottomrule
\end{tabular}
\end{table}

\subsubsection*{表关系与隐私设计原则}

四张表之间通过外键关联，整体以订单号为核心索引构建数据流。\texttt{users} 与 \texttt{orders} 通过 \texttt{user\_id} 关联，\texttt{orders} 与 \texttt{messages} 通过 \texttt{order\_id} 关联，\texttt{audit\_logs} 通过 \texttt{order\_id} 与订单和消息建立审计关系。

在隐私保护方面，数据库设计遵循以下原则：其一，消息明文不落库，\texttt{content}字段存储 SM4-CBC 密文，明文只在服务端内存中短暂存在；其二，\texttt{pid\_r}（随机数$r$）与\texttt{pid}均仅服务端维护，前端不可见，骑手端不可见；其三，\texttt{note}备注字段在 API 响应中对骑手端和管理员端屏蔽，仅在用户本人查询订单时返回；其四，\texttt{sender\_pid}字段中存储的是 PID 而非用户名，管理员默认视图也不直接暴露 PID 对应的真实用户名。
